\chapter*{Introducción}
\addcontentsline{toc}{chapter}{Introducción}

La Holografía Modular Triádica parte de una decisión elemental y de gran
alcance: conservar, junto al resultado de una operación, la información que
la operación ordinaria descarta. Al dividir un entero por nueve aparecen un
residuo y un cociente. La aritmética usual retiene el valor recompuesto;
HMT retiene además el observable de procedencia, la orientación, el acarreo
fronterizo y la clase de prolongaciones compatibles. De esa decisión nace una
aritmética de estados provistos de un registro de transporte.

Su objeto inicial es la estructura de todos los caminos posibles
(\emph{All Possible Paths}, APP), denotada por \(\mathcal A_9\). Sobre un
mismo soporte finito, APP reúne los observables de suma y producto, las rutas
orientadas y la extensión residuo--cociente. El estado ternario orientado se
denomina TRIT, y el sistema dinámico que transporta orientación, frontera y
fase es el núcleo topológico de fase (\emph{Topological Phase Kernel}, TPK).
Estas tres piezas organizan un catálogo exhaustivo de historias finitas. El
catálogo puede proyectarse sobre palabras ternarias, extenderse mediante
coordenadas de cociente módulo nueve y completarse mediante cilindros encajados.

La construcción matemática tiene dos direcciones coordinadas. En la
primera, las historias compatibles se contraen mediante una evaluación
arquimediana: un punto real aparece como límite de una familia de cilindros.
En la segunda, una frontera marcada, junto con la calibración de su codominio,
se despliega como soportes, códigos, diseños, retículos y automorfismos
excepcionales. El valor procede de la evaluación arquimediana; la simetría,
de la prolongación combinatoria. El teorema finito de compatibilidad muestra
que ambas operaciones producen la bandera \(B_K\subset H^-_\alpha\) para los
datos dodecafásicos y combinatorios declarados. La extensión global se
formula mediante los dominios y mapas tipados del capítulo correspondiente.

El desarrollo se realiza desde lo finito hacia lo global. Primero se prueba
la contabilidad exacta del soporte APP, incluida la identidad
\[
2025-810=(459-405)+9(174-45)=54+9\cdot129.
\]
Después se construye la aplicación de codificación asociada a dos recorridos
posicionales y se enumera todo su dominio:
\[
104\,976\longrightarrow468\longrightarrow243.
\]
La imagen del codificador contiene exactamente 42 ternas decimales y satisface
una congruencia de control; las 468 salidas se factorizan en 26 firmas y 18 representantes de
la acción de fase. La fibra asociada a \(\pi\) posee cinco regiones
distintas, y las cinco contienen paseos cerrados mínimos de
longitud ocho junto a los anclajes de \(e\), \(\varphi\) y APP\@.

En el nivel dodecafásico, fijado como premisa enriquecida el vector firmado
\(U_{12}^{\rm sgn}\), una transformación reversible determina el vector
\(K\). Tres representaciones reversibles producen la reconstrucción
interna de \(\alpha\), y el teorema correspondiente demuestra la composición
exacta y la conmutatividad del diagrama que las relaciona.

La rama excepcional comienza en una matriz Paley--Witt sobre
\(\mathbb F_3\), construye el código ternario extendido, el diseño
\(W_{12}\) y su grupo \(M_{12}\). Una rama binaria hacia \(W_{24}\) se
define mediante su mapa de incidencia. La otra fija el pegado del Niemeier
\(A_2^{12}\), selecciona un 3-vecino y certifica que es par, unimodular y sin
raíces. La clasificación lo identifica con Leech, desde donde los teoremas
clásicos construyen \(V^\natural\) y el grupo Monstruo. Las dos ramas se
mantienen tipadas y separadas.

Finalmente se construye una multiplicación global a partir de las operaciones
locales. El producto nativo
\(\mathcal A_9^\#\) se construye celda a celda con las coordenadas de residuo y
cociente, se normaliza por el observable aditivo y se extiende por Horner a palabras
canónicas. El entrelazador con los enteros aparece entonces como teorema de
conservación, no como definición circular.

Cada afirmación finita importante posee en esta carpeta sus datos y su
procedimiento computacional de reproducción. Las pruebas Lean incluidas verifican un conjunto
inicial de identidades algebraicas con el núcleo lógico del asistente. El
texto impreso contiene las demostraciones conceptuales; los programas
enumeran dominios completos y suministran testigos; los certificados fijan
las salidas. Las tres formas están diseñadas para poder leerse y comprobarse
sin recurrir a archivos situados fuera de esta obra.
